\documentclass[12pt]{article}

% page layout
\usepackage[a4paper, margin=1in]{geometry}
\usepackage{lineno}
\usepackage{setspace}

% misc
\usepackage{multirow}
\usepackage{csquotes}

\usepackage{microtype}
% \overfullrule=5pt

% graphics and colors
\usepackage[dvipsnames]{xcolor}
\usepackage{graphicx}
\usepackage{float}

% small font size for captions
\usepackage[font=footnotesize,labelfont=bf]{caption}

% math and tables
\usepackage{amsmath}
\usepackage{booktabs}
\usepackage{siunitx}
\sisetup{retain-explicit-plus}
\DeclareSIUnit{\eggs}{eggs}
\DeclareSIUnit{\year}{year}
\DeclareSIUnit{\Otwo}{O\textsubscript{2}}

% bibliography (Ecology Letters style)
\usepackage[
  style=authoryear,
  backend=biber,
  sorting=nyt,
  sortcites=true,
  maxcitenames=2,
  mincitenames=1,
  maxbibnames=6,
  minbibnames=6,
  uniquename=false,
  uniquelist=false,
  giveninits=true,
  dashed=false,
  date=year,
  labeldate=year,
  doi=false,
  url=false,
  isbn=false,
  eprint=false
]{biblatex}
% citations: "Author Year", "&" between two authors
\DeclareDelimFormat{nameyeardelim}{\addspace}
\DeclareDelimFormat{finalnamedelim}{\addspace\&\space}
\DeclareDelimFormat[bib]{finalnamedelim}{\addspace\&\space}
\DeclareNameAlias{sortname}{family-given}
% reference list: "Title. Abbrev. J., Volume, pages."
\DeclareSourcemap{
  \maps[datatype=bibtex]{
    \map[overwrite]{
      \step[fieldsource=shortjournal]
      \step[fieldset=journaltitle, origfieldval]
    }
    \map{
      \pertype{article}
      \step[fieldset=editor, null]
    }
  }
}
\DeclareFieldFormat[article]{title}{#1}
\DeclareFieldFormat{journaltitle}{#1\isdot}
\DeclareFieldFormat{pages}{#1}
\renewbibmacro{in:}{\ifentrytype{article}{}{\printtext{\bibstring{in}\intitlepunct}}}
\renewbibmacro*{journal+issuetitle}{%
  \usebibmacro{journal}%
  \setunit{\addcomma\space}%
  \printfield{volume}%
  \setunit*{\addcomma\space}%
  \printfield{eid}}

% authors and affiliations
\usepackage{authblk}
\renewcommand*{\Affilfont}{\footnotesize\normalfont}

% listings and text code formatting packages
\usepackage{listings}

% hyperlinks
\usepackage{hyperref}
\hypersetup{
  colorlinks=true,
  linkcolor=Blue,
  filecolor=Maroon,
  citecolor=Blue,
  urlcolor=Blue
}

% custom commands for title page
\title{Beyond the average phenotype: individual variation alters mechanistic forecasts of climate responses}

\author[1*]{Urtzi Enriquez-Urzelai}
\author[1,2]{Gioele Pascucci}
\author[1]{Lumír Gvoždík}

\affil[1]{Czech Academy of Sciences, Institute of Vertebrate Biology, Brno, Czech Republic}
\affil[2]{Department of Botany and Zoology, Faculty of Science, Masaryk University, Brno, Czech Republic}
\affil[ ]{E-mails: UEU, uenriquez@proton.me; GP, gioelepascucci@mail.muni.cz; LG, gvozdik@ivb.cz}
\affil[*]{Corresponding author}
\date{\vspace{-5ex}}

\makeatletter
% fix authblk spacing
\renewcommand\AB@affilsepx{\protect\\\vspace{\baselineskip}\protect\Affilfont}

% store title and author for appendix
\let\papertitle\@title
\let\paperauthors\@author

% appendix title block
\newcommand{\appendixtitleblock}[1]{
  \begin{center}
    {\Large\textbf{Appendix #1} \textit{for:} \papertitle}
    \vskip 1em
      {\large\paperauthors}
  \end{center}
  \vskip 1.5em
}

% image width constraint
\def\maxwidth{%
  \ifdim\Gin@nat@width>\textwidth
    \textwidth
  \else
    \Gin@nat@width
  \fi
}
\makeatother

% custom counter for appendix cross-references
\newcounter{appnum}

% custom macro to setup the appendices
\newcommand{\appendixsetup}[1]{
  \setcounter{appnum}{\numexpr#1-1\relax}%
  \refstepcounter{appnum}
  \newpage
  \singlespacing
  \ifdefined\nolinenumbers\nolinenumbers\fi
  \setcounter{page}{1}
  \newrefsection
  \setcounter{section}{0}
  \renewcommand{\theHsection}{app#1.\arabic{section}}
  \renewcommand{\thefigure}{S#1.\arabic{figure}}
  \renewcommand{\theHfigure}{S#1.\arabic{figure}}
  \setcounter{figure}{0}
  \renewcommand{\thetable}{S#1.\arabic{table}}
  \renewcommand{\theHtable}{S#1.\arabic{table}}
  \setcounter{table}{0}
}
\addbibresource{references.bib}
\addbibresource{packages.bib}

\makeatletter
\def\maxwidth{ %
  \ifdim\Gin@nat@width>\linewidth
    \linewidth
  \else
    \Gin@nat@width
  \fi
}
\makeatother

\begin{document}



\doublespacing
\maketitle
\thispagestyle{plain}

\noindent\textbf{Running title:} Individual variation alters forecasts

\noindent\textbf{Keywords:} biophysical modelling; climate change; Dynamic Energy Budget theory; ecological forecasting; individual variation; intraspecific variation; Jensen's inequality; nonlinear averaging; performance surface

\noindent\textbf{Type of article:} Letter

\noindent\textbf{Word counts:} abstract, 147; main text, 3945; text boxes, none

\noindent\textbf{Number of references:} 73

\noindent\textbf{Number of figures, tables and text boxes:} 4 figures; no tables; no text boxes

\noindent\textbf{Correspondence:} Urtzi Enriquez-Urzelai, Institute of Vertebrate Biology, Czech Academy of Sciences, Brno, Czech Republic. E-mail: uenriquez@proton.me

\noindent\textbf{Statement of authorship:} UEU and LG designed the study. LG provided the individual growth and metabolic data. UEU estimated the DEB parameters, performed the simulations and analyses, and produced the figures. UEU wrote the first draft of the manuscript, and all authors contributed substantially to revisions.

\noindent\textbf{Data accessibility statement:} All data and code supporting the results are archived in Zenodo (\url{https://doi.org/10.5281/zenodo.XXXXXXX}; \cite{enriquez-urzelai2026data}). The individual growth and metabolic data were originally published by \textcite{baskiera2021}.

\newpage
\linenumbers

\section*{Abstract}
\noindent
Individuals differ in physiology, and organismal performance depends nonlinearly on both phenotype and environment. Jensen's inequality therefore predicts that the performance of an average individual differs from the average performance of a population, yet how this bias varies across traits and environments has rarely been tested with measured individual variation. We estimated Dynamic Energy Budget parameters for 37 individually monitored alpine newts (\emph{Ichthyosaura alpestris}) and for their population average, and coupled them with biophysical models to predict twelve performance traits across Europe under current climate and \qty{+2}{\celsius} and \qty{+4}{\celsius} warming. Average-phenotype predictions diverged systematically from mean individual predictions, with the sign and magnitude of this Jensen bias depending on trait, sex and location; reproduction was the most sensitive. Warming amplified, attenuated or reversed the bias, sometimes non-monotonically, although both approaches agreed on the direction of change. The consequences of individual variation are therefore strongly context dependent.

\newpage

\section{Introduction}

Ecologists routinely predict population-level outcomes from mean trait values. Trait-based and mechanistic niche models, for instance, commonly summarise a population or species by an average phenotype \parencite{kearney2009,laughlin2015,briscoe2023}. Yet individuals within populations consistently differ in morphology, physiology, behaviour and life history, and these differences influence resource use, species interactions and demographic rates \parencite{bolnick2003,bolnick2011,desroches2018}. Whether an average phenotype adequately represents a variable population depends on how trait variation is translated into performance. When this translation is nonlinear, the performance of the average individual differs from the average performance of individuals, a consequence of Jensen's inequality \parencite{jensen1906,ruel1999}. Although this principle is well established in theory, how strongly it biases ecological predictions, and how that bias changes across traits and environments, has rarely been quantified with empirically measured individual variation.

Performance emerges from interactions between phenotype and environment. These interactions can be represented as performance surfaces describing how a given component of performance varies across combinations of phenotypes and abiotic conditions (\cite{laughlin2015,li2021}; Fig.~\ref{fig:conceptual}a). Such surfaces are rarely linear: thermal optima, thresholds and asymmetries are pervasive in ecological systems \parencite{ruel1999,bolnick2011}. We define Jensen bias as $J = f(\overline{\theta}) - \overline{f(\theta)}$, the difference between performance predicted from the mean phenotype $\overline{\theta}$ and mean performance across phenotypes. To second order, $J \approx -\frac{1}{2} f''(\overline{\theta})\operatorname{Var}[\theta]$, so bias arises not from trait variation alone but from its interaction with the local curvature of the performance surface \parencite{denny2017,denny2019}. Where the surface is concave, the average phenotype overestimates mean performance ($J>0$); where it is convex, it underestimates it ($J<0$; \cite{ruel1999,bolnick2011}). This simple result has three consequences that are rarely considered together (Fig.~\ref{fig:conceptual}). First, because curvature changes along environmental gradients, the same distribution of phenotypes can produce negligible bias in some environments and substantial bias in others \parencite{ross2023}. Second, because growth, reproduction and survival need not share the same optima or thresholds \parencite{sinclair2016,kontopoulos2024}, different components of performance can be biased by different amounts and even in opposite directions. Third, any environmental change that moves populations to a different region of the surface can amplify, attenuate or reverse the bias (Fig.~\ref{fig:conceptual}b,c).

Physiology provides a natural setting in which to test these predictions, because it is the mechanistic link between environmental conditions and ecological performance \parencite{seebacher2012}, and individuals commonly differ in metabolic rates, thermal sensitivity and energy allocation \parencite{pettersen2024,gvozdik2024}. Mechanistic models make this link explicit by simulating how organisms interact with their microclimates to determine growth, survival and reproduction \parencite{kearney2009,kearney2020c,briscoe2023}. These models are increasingly used to forecast responses to climate change and now represent environmental heterogeneity in great detail, yet they almost always represent species by a single average phenotype (but see \cite{kolbe2010,enriquez-urzelai2020,riddell2022}). Because their performance surfaces are generated from first principles rather than assumed, mechanistic models allow measured individual variation to be propagated through realistic nonlinearities, while climate warming provides an environmental shift of known magnitude with which to test whether the bias changes as environments change.

Here, we combined longitudinal measurements of growth and metabolism from individually monitored alpine newts (\emph{Ichthyosaura alpestris}; data from \cite{baskiera2021}) with Dynamic Energy Budget (DEB) and biophysical models to compare predictions from a population-average parameterisation with the mean of predictions from 37 individual parameterisations, under identical environmental simulations across Europe and under warming. We tested three hypotheses: (H1) predictions from the average phenotype diverge systematically from the mean of individual predictions; (H2) the sign and magnitude of this Jensen bias differ among performance components and across environmental gradients, reflecting differences in local performance-surface geometry; and (H3) warming changes Jensen bias in trait- and context-specific ways, including amplification, attenuation and reversal. We found support for all three hypotheses, although the two parameterisations largely agreed on the direction of the projected climate response.

\section{Materials and Methods}

\subsection{Study system and data}

The alpine newt is a temperate amphibian widely distributed across Europe. Juveniles are strictly terrestrial until sexual maturity, and adults spend most of the year on land, outside the breeding season (April--June), where they are predominantly nocturnal or crepuscular \parencite{denoel2004,rocek2003}. We compiled species-level life-history and growth information from the literature, including lifespan (20 years), snout--vent length (SVL) at metamorphosis (\qty{1.83}{\cm}), female and male SVL at sexual maturity (4.9 and \qty{4.4}{\cm}), asymptotic SVL (\qty{5.84}{\cm} females, \qty{5.05}{\cm} males), asymptotic body mass (\qty{3.16}{\g} females, \qty{2.04}{\g} males), maximum reproductive rate (\qty{600}{\eggs\per\year}) and published growth trajectories \parencite{gvozdik2017,miaud2000a}. Individual-level data consisted of annual measurements of body mass (\si{\gram}) and standard metabolic rate (\si{\milli\liter\Otwo\per\hour}) of 37 alpine newts (18 females and 19 males) from a population near Jihlava, Czech Republic (\ang{49.2247}\,N, \ang{16.0579}\,E; \qty{600}{m} a.s.l.), monitored from metamorphosis in 2014 until 2018 \parencite{baskiera2021}. These data were collected in a previous study under the permits and ethical approvals reported therein \parencite{baskiera2021}; no new animals were sampled for this study.

\subsection{DEB parameterisation}

DEB theory describes the acquisition, storage and allocation of energy throughout the life cycle \parencite{kooijman2010}. The standard DEB model represents an organism by its structure, reserve, maturity and, in adults, a reproduction buffer, and is governed by 12 primary parameters, each quantifying the rate or efficiency of a specific physiological process \parencite{kooijman2010,marques2018}. We estimated parameters by inverse modelling using the covariation method \parencite{lika2011,marques2018}, with the AmPtool and DEBtool\_M packages in MATLAB (The MathWorks Inc., Natick, MA, USA), at two levels (Appendix~S1). First, a single population-average parameterisation (hereafter, average phenotype) was fitted simultaneously to the species-level information and to all longitudinal observations of the 37 individuals, retained as separate data points. This mirrors standard practice in mechanistic modelling and avoids averaging nonlinear trajectories into a composite that no individual follows. We estimated seven primary parameters, namely the zoom factor $z$ (which scales maximum size and surface-specific assimilation; \cite{sousa2008}), energy conductance $\dot{v}$, allocation fraction $\kappa$, volume-specific somatic maintenance $[\dot{p}_M]$, specific cost of structure $[E_G]$ and the maturity thresholds at birth $E_H^b$ and puberty $E_H^p$, together with a threshold at metamorphosis and a male-specific zoom factor ($z_m$) and puberty threshold ($E_{H,m}^p$); the remaining parameters were fixed at standard values \parencite{kooijman2010,kearney2013}.

Second, we re-estimated three parameters independently for each individual from its own growth and metabolic trajectories: the sex-specific zoom factor ($z$ or $z_m$), $\dot{v}$ and $[\dot{p}_M]$. These parameters govern the processes captured by the individual data, i.e., energy uptake and body size, reserve mobilisation and maintenance costs \parencite{kooijman2021}. All other parameters were held at their population-average values, and all individuals were assigned the same temperature (\qty{12}{\celsius}) and food availability during estimation, so that differences among parameter sets reflect differences in growth and metabolism rather than in the conditions experienced. Parameter values are given in Table~S2.1.

\subsection{Mechanistic simulations}

We coupled each DEB parameterisation with biophysical models using the NicheMapR R package \parencite{kearney2017a,kearney2020c}. The biophysical model computes hourly body temperature, activity and foraging time from microclimatic conditions and the organism's morphology, physiology and behaviour \parencite{porter1973,kearney2018}. Body temperature then scales all DEB rates through the Arrhenius relationship and, with food availability held constant, foraging time determines energy intake \parencite{kearney2012a,kearney2020c}. Hourly microclimates were computed with \texttt{micro\_global()} from the 10~arc-minute climatology of \textcite{new2002}, for 0\% and 90\% shade, with moisture and snow routines enabled, for 22 years, of which the first two were discarded as spin-up \parencite{kearney2014,kearney2018a,kearney2020b}.

The \texttt{ectotherm()} model was parameterised for a terrestrial newt with 90\% of the skin acting as a free water exchanger \parencite{enriquez-urzelai2024}. Animals were active at night and at dusk (juveniles also by day; \cite{enriquez-urzelai2020}) when body temperature was between \qty{4}{\celsius} and \qty{20}{\celsius} and relative humidity was at least 65\% \parencite{bonnefond2020,giacometti2025}, selecting shade (0--90\%) to approach the preferred body temperature of \qty{17.3}{\celsius} \parencite{gvozdik2017}; otherwise they retreated underground to depths (\qty{2.5}{\cm}--\qty{2}{\m}) that kept them within critical thermal limits (\qty{-2}{\celsius} and \qty{36}{\celsius}; \cite{enriquez-urzelai2024}). Females reproduced in batches of 200--600 eggs \parencite{pecquerie2009}, whereas males depleted their reproduction buffer continuously during the breeding season. Simulations started with juveniles on 15 August and ran until death or for up to 20 years, the maximum known lifespan (Appendix~S1).

Simulations were run for every terrestrial cell of a 10~arc-minute grid covering Europe (\ang{10}\,W--\ang{40}\,E, \ang{35}--\ang{70}\,N), for the female and male average phenotypes and for each of the 37 individual parameterisations. The whole framework was repeated under uniform increases of \qty{+2}{\celsius} and \qty{+4}{\celsius} in hourly air temperature, which span projected warming across Europe under intermediate to high emission pathways \parencite{calvin2023}. All other climatic variables were held constant, isolating the thermal effects of warming; because relative humidity was unchanged, warming also increased evaporative demand. Thus, within each grid cell and scenario, predictions differed only in the physiological parameterisation.

\subsection{Quantifying Jensen bias}

For every grid cell, parameterisation and scenario, we extracted 12 performance metrics: lifespan (\si{\year}), lifetime reproductive output (eggs in females, hereafter lifetime fecundity; reproductive energy, \si{\joule}, in males), net reproductive rate ($R_0$), maximum and minimum annual body temperature ($T_{b,max}$, $T_{b,min}$), maximum body mass (\si{\g}), maximum SVL (\si{\cm}), minimum reserve energy (\si{\joule}), and lifetime basking time (\si{\hour}; time above ground without foraging), foraging time (\si{\hour}), food intake (\si{\g} dry mass) and mean daily water need (g~day$^{-1}$). For each metric and sex, Jensen bias was computed per grid cell as
\[
J=P_{\mathrm{pop}}-\overline{P}_{\mathrm{ind}},
\]
where $P_{\mathrm{pop}}$ is the prediction of the average phenotype and $\overline{P}_{\mathrm{ind}}$ the mean of the individual predictions. Positive values indicate overprediction by the average phenotype and negative values underprediction. To compare metrics with different units, we used the normalised absolute bias: the mean $|J|$ across grid cells divided by the mean $|\overline{P}_{\mathrm{ind}}|$ ($\times 100$). Independently of bias, we quantified among-individual divergence in predictions as the coefficient of variation (CV) of individual predictions in each cell, and compared its median between southern ($<$\ang{40}\,N), central and northern ($>$\ang{60}\,N) latitudes.

To test how warming modified bias, we computed $\Delta |J|=|J_{\mathrm{future}}|-|J_{\mathrm{current}}|$, where positive values indicate amplification and negative values attenuation, and summarised it per metric as the relative change, i.e., the mean $\Delta|J|$ divided by the mean $|J_{\mathrm{current}}|$ ($\times 100$). We also report the proportions of cells in which $|J|$ increased or decreased and in which the sign of $J$ reversed. We described how $\Delta|J|$ under \qty{+4}{\celsius} varied along mean annual air temperature ($\overline{T}_a$) using generalised additive models fitted with \texttt{mgcv} \parencite{R-mgcv}; because model outputs are deterministic and spatially autocorrelated, these served as descriptive smoothers rather than for inference (Appendix~S1).

Finally, for three traits summarising survival, growth and reproduction (lifespan, maximum body mass and lifetime fecundity or reproductive energy), we compared predictions of both parameterisations under each scenario and their projected changes ($P_{\mathrm{future}} - P_{\mathrm{current}}$). We report the percentage of cells in which the average phenotype under- or overpredicted the individual mean, the percentage in which both projected the same direction of change, and the median and 95th percentile of absolute relative bias, $(P_{\mathrm{pop}} - \overline{P}_{\mathrm{ind}}) / |\overline{P}_{\mathrm{ind}}| \times 100$. All data and code are available in Zenodo \parencite{enriquez-urzelai2026data}.

\section{Results}

\subsection{Individual physiological variation}

Individual DEB parameters varied among the 37 individuals, with coefficients of variation ranging from 1.7\% ($z$) to 8\% ($[\dot{p}_M]$; Fig.~S2.1). Somatic maintenance $[\dot{p}_M]$ was the most variable (CV = 8\%), followed by energy conductance $\dot{v}$ (CV = 5.6\%) and the zoom factors ($z$: 1.7\%; $z_m$: 1.9\%), and each parameter was approximately normally distributed around the population-average estimate (Fig.~S2.2). Both parameterisations reproduced the observed growth and metabolic trajectories accurately (Figs.~S2.3, S2.4, S2.5--S2.14).

\subsection{Physiological averaging generates Jensen bias (H1, H2)}

Under the current climate, predictions of the average phenotype differed systematically from the mean of individual predictions for most metrics. For females, the average phenotype underestimated most metrics (Fig.~\ref{fig:map-jensen}), including basking time ($|J| = 419$\,h), foraging time ($|J| = 143$\,h), lifetime fecundity ($|J| = 63$\,eggs), food intake ($|J| = 3.7$\,\si{\g}) and maximum body mass ($|J| = 0.2$\,\si{\g}; absolute value of the mean $J$ across cells). Relative to the individual mean, bias was largest for $R_0$ (normalised $|J| = 13.3$\%), lifetime fecundity (10.1\%) and maximum body mass (5.3\%), and negligible for body temperatures ($\leq 0.2$\%; mean $|J| < 0.004$\,${}^{\circ}$C). For males, the pattern was largely opposite: the average phenotype overestimated most metrics, with the largest relative bias for reproductive energy (normalised $|J| = 9.1$\%) and $R_0$ (7.4\%; Fig.~S3.1, Appendix~S3).

The direction of bias differed among metrics. In females, underestimation dominated across most of the range for body size, reserve, food intake, lifetime fecundity and basking ($J<0$ in 63--98\% of cells), whereas body temperatures were predominantly overestimated ($J>0$ in 65--83\% of cells), and foraging time was split between overestimation (52\% of cells), underestimation (46\%) and no bias (2\%). In males, these directions were reversed ($J>0$ in 57--96\% of cells for the same metrics).

Among-individual divergence in predictions was also geographically structured (Figs.~S3.2, S3.3). Compared with central latitudes, the median CV of individual predictions in females was 3.3-fold higher at the southern margin for lifetime fecundity, 3.1-fold for minimum reserve and 3.4-fold for maximum body length. In males, the largest increases occurred at the northern margin for $R_0$ (4.3-fold) and reproductive energy (2.2-fold). In contrast, lifespan, basking time, foraging time and body temperatures differed by at most 1.5-fold between marginal and central latitudes in either sex.

\subsection{Climate warming reshapes Jensen bias (H3)}

Under \qty{+4}{\celsius} warming, $|J|$ changed for all 12 female metrics (Fig.~\ref{fig:map-jensen-warming}). The strongest relative amplification occurred for minimum body temperature (rel.\ $\Delta|J| = +$568\%; increase in 72\% of cells), maximum body temperature (rel.\ $\Delta|J| = +$182\%; 73\% of cells) and $R_0$ (rel.\ $\Delta|J| = +$138\%), whereas lifespan ($-15\%$) and minimum reserve ($-6\%$) were attenuated. The large relative changes in body temperature bias, however, arose from very small absolute biases (mean $|J|$ from 0.0021 to 0.014\,${}^{\circ}$C for $T_{b,min}$ and from 0.0035 to 0.0099\,${}^{\circ}$C for $T_{b,max}$). In males, the strongest amplification occurred for minimum body temperature (rel.\ $\Delta|J| = +$510\%) and maximum body temperature (rel.\ $\Delta|J| = +$215\%), whereas lifespan ($-1\%$) was attenuated (Fig.~S3.4).

Warming also reversed the direction of bias in some regions. The sign of $J$ changed between current and \qty{+4}{\celsius} scenarios in 2--49\% of cells in females, most often for minimum body temperature (49\%), foraging time (47\%) and $R_0$ (44\%), and in 5--46\% of cells in males. Changes in bias were geographically structured and varied systematically, but in metric-specific ways, along the mean annual temperature gradient (Fig.~\ref{fig:map-jensen-warming}). Responses were not always monotonic: under \qty{+2}{\celsius}, 11 female metrics showed increased $|J|$ (from $+$4\% for daily water need to $+$152\% for minimum body temperature), and lifespan switched from amplification at \qty{+2}{\celsius} to attenuation at \qty{+4}{\celsius} (see Appendix~S3 for males).

\subsection{Consequences for forecasts}

The average phenotype underestimated female and overestimated male performance for several traits, with differences in the magnitude and geographical extent of bias (Figs.~\ref{fig:map-suitability}, S3.5). In females, underprediction of lifetime fecundity became more widespread with warming, from 63\% of cells under the current climate (median relative bias $-7.1$\%) to 71\% and 76\% under \qty{+2}{\celsius} and \qty{+4}{\celsius}. Underpredicted lifespan increased from 76\% of cells (median $-1.9$\%) to 86\% under \qty{+4}{\celsius}, while maximum body mass was underpredicted in 83\% of cells (median $-3.8$\%). Across 95\% of cells, the average phenotype underestimated female lifespan by up to 19\% and lifetime fecundity by up to 100\%, compared with $\pm$19\% for body mass (Fig.~S3.6). In males, reproductive energy was overpredicted in 85\% of cells under the current climate (median $+8.1$\%) and 95\% under \qty{+4}{\celsius}, by up to 66\% of the individual mean (Fig.~S3.7). Despite these differences, both parameterisations projected the same direction of change under \qty{+4}{\celsius} in at least 90\% of cells for all three traits in both sexes.

\section{Discussion}

Jensen's inequality implies that the ecological consequences of individual variation depend on the shape of the relationship between phenotype, environment and performance \parencite{ruel1999,bolnick2011}. By propagating measured physiological variation through mechanistic performance surfaces, we show that this principle has large, structured consequences for predictions: an average phenotype misrepresented the mean performance of a heterogeneous population for most traits (H1), the bias differed in magnitude and sign among traits, sexes and locations (H2), and warming reshaped it in trait-specific and sometimes non-monotonic ways (H3). Notably, variation in only three energetic parameters was sufficient to generate these patterns at continental scale, extending previous work on nonlinear responses to environmental variability \parencite{paaijmans2013,denny2017} to variability among individuals.

\subsection*{Which traits are sensitive to averaging}

Reproductive traits showed the largest relative bias in both sexes, whereas body mass was less sensitive. This difference is consistent with how traits integrate energy over the life cycle. Reproduction reflects acquisition, allocation and maturation history throughout adulthood \parencite{kooijman2010,jusup2017}, so small physiological differences accumulate through time and are amplified by thresholds such as the energy required for a clutch. Body mass, in contrast, is constrained by tight feedbacks between assimilation and maintenance that converge on an asymptotic size \parencite{nisbet2012,lika2019}, buffering the effect of individual variation. Traits that integrate cumulative allocation are therefore expected to be more sensitive to averaging than structural traits \parencite{nisbet2012,palmer2024}, a prediction that should apply beyond our study system. The opposite direction of bias in females and males further shows that the sign of $J$ is not a property of the trait alone: batch reproduction in females converts energetic differences into discrete clutches, whereas continuous depletion of the reproduction buffer in males smooths them \parencite{pecquerie2009,muller2019}, so that the same parameter variation meets performance surfaces of different curvature. Likewise, foraging time switched between over- and underestimation across the range, indicating that local environmental conditions can determine the direction of bias \parencite{martin2008}.

\subsection*{Where averaging matters}

The divergence among individual predictions increased towards the northern and southern margins of the modelled area, particularly for reproduction- and reserve-related traits. This is consistent with the expectation that the consequences of trait variation depend on where populations operate on the performance surface, which changes along environmental gradients \parencite{laughlin2015,li2021,ross2023,polazzo2024}. Under benign conditions, individuals with different physiologies may achieve similar outcomes, whereas under constraining conditions small differences can determine whether energy suffices for growth, reserves or reproduction. Although the strongest effects occurred partly beyond the current range of \emph{I.~alpestris} \parencite{iucn2021}, they emerged precisely where ecological responses are expected to become most sensitive, as populations approach environmental limits \parencite{holt2003,sexton2009,valladares2014}. Predictions based on average phenotypes are therefore least reliable where predicting persistence, range expansion and local extinction matters most \parencite{demarche2019}.

\subsection*{Environmental change reshapes the bias}

Warming changed Jensen bias for every metric, but not uniformly. It amplified the bias in net reproductive rate in both sexes, the most demographically relevant change, attenuated it for lifespan, and reversed its sign across large parts of the range for some metrics. For lifespan, bias increased under \qty{+2}{\celsius} but decreased under \qty{+4}{\celsius}, showing that the effect of individual variation does not scale monotonically with environmental change. These responses are expected when warming moves populations across regions of a nonlinear surface with different curvature \parencite{angillettajr.2009,kingsolver2016,li2021}: further warming can move a population past an inflection point, weakening rather than strengthening the effect of averaging. Hence, whether an average phenotype over- or underestimates performance is a property of the phenotype distribution, the trait and the environment jointly, and validation of average-phenotype predictions under current conditions does not guarantee their accuracy under future conditions \parencite{buckley2018,buckley2023,cannistra2024}.

\subsection*{Robust trends, biased magnitudes}

Despite these biases, both parameterisations agreed on the direction of the projected response in at least 90\% of cells: warming reduced lifespan across most of the range \parencite{munch2009,burraco2020}, and both projected reduced body mass and reproduction towards the north-eastern range limit, where warming improves thermal suitability. Physiological averaging therefore affected quantitative predictions, i.e., the magnitude and spatial pattern of change, without overturning the qualitative interpretation. This distinction between robust trends and biased magnitudes clarifies when average phenotypes are adequate: for the direction of change, our results suggest they generally are, whereas predictions of absolute performance, and particularly of reproduction and population growth, should incorporate individual variation. Mechanistic models are well suited to this task because they represent physiology explicitly and have already been used to explore trait evolution and thermal adaptation using observed or hypothetical variation \parencite{dejong1996,kingsolver2016,enriquez-urzelai2020,riddell2022,riddell2023a}. Propagating empirical trait distributions requires that individual-level estimates, not only population means, are reported and archived \parencite{wolf2023}.

Several simplifications point to future work. We allowed only three parameters to vary, used individuals from a single population and held trait distributions fixed under warming, whereas variation in other traits, among-population differentiation and plasticity or adaptation are all expected to shape trait distributions across ranges and through time \parencite{moran2016,sexton2009,chevin2010,kingsolver2017}. Each of these could amplify, attenuate or reverse the patterns of Jensen bias reported here, and integrating them is an important next step towards biologically realistic predictions.

In conclusion, representing a population by a single average phenotype introduced systematic, spatially structured and climate-dependent biases in predicted performance, because individual variation interacts with nonlinear phenotype--environment--performance relationships. The bias was largest for traits that integrate energy allocation over the life cycle and under environmentally constraining conditions, and it could be amplified, attenuated or reversed by environmental change. More generally, our results show that the ecological consequences of individual variation cannot be inferred from the amount of variation alone, but depend on where populations lie on the underlying performance surface, which is a principle that applies to any prediction made from mean traits in a changing environment.

\section*{Acknowledgements}

We thank Bas Kooijman and Gonçalo Marques for help with DEB parameter estimation, and Michael R. Kearney for help setting up the sex-specific NicheMapR simulations. Computational resources were provided by the e-INFRA CZ project (ID:90254), supported by the Ministry of Education, Youth and Sports of the Czech Republic. Funding was provided by the Czech Science Foundation (25-15551S to UEU, GP and LG), the Department of Botany and Zoology, Faculty of Science, Masaryk University in Brno, Czechia (MUNI/A/1615/2025 to GP), and the Institute of Vertebrate Biology of the Czech Academy of Sciences (RVO: 68081766 to LG). The authors declare no conflicts of interest.

% REFERENCES section
\printbibliography[heading=bibnumbered, title={References}]

% FIGURES (grouped at the end, as required by Ecology Letters)
\clearpage
\nolinenumbers

\begin{figure}[p]
  \centering
  \includegraphics[width=\textwidth]{/home/urtzi/Documents/papers/individuals_DEB/figures/2_final/fig1-conceptual.pdf}
  \caption{Jensen's inequality and its implications for predicting performance in changing environments. For a nonlinear function, the value at the mean $f(\overline{\theta})$ differs from the mean of the values $\overline{f(\theta)}$, producing a bias (Jensen bias, $J = f(\overline{\theta}) - \overline{f(\theta)}$), where $J \approx -\frac{1}{2} f''(\overline{\theta})\,\operatorname{Var}[\theta]$, whose sign and magnitude follow the local curvature of $f$ \parencite{denny2017,jensen1906}.
  \textbf{(a) Representative trait--performance surfaces illustrating common forms of nonlinear responses to temperature}. The blue-to-red bar on the far left represents an environmental temperature gradient (cold to warm). Performance ($F$) is shown as a function of trait value ($\theta$) and temperature ($T$) for five illustrative types: concave dome, Gaussian, Gaussian with shifting thermal optimum ($T_{opt}$), sigmoidal with fixed threshold, and sigmoidal with shifting threshold. Coloured vertical slices show the trait--performance relationship at cold (\emph{blue}), intermediate (\emph{black}) and warm (\emph{red}) temperatures, while \emph{black points} denote the mean trait value.
  \textbf{(b) Temperature-dependent Jensen bias predicted for each response surface}. \emph{Red shading} indicates positive bias ($J > 0$), where predictions based on the mean trait overestimate performance, whereas \emph{blue shading} indicates negative bias ($J < 0$), where the mean trait underestimates performance. The \emph{solid curve} represents the bias under the current climate, $J(T)$, and the \emph{dashed curve} the bias after warming, $J(T+\Delta T)$, for a fixed trait distribution ($\bar\theta$ and $\operatorname{Var}[\theta]$ held constant).
  \textbf{(c) Impact of warming on the magnitude of Jensen bias}. The \emph{black curve} shows the change in absolute bias, $|J|$, across the temperature gradient. \emph{Yellow shading} denotes amplification ($|J|$ increases), whereas \emph{purple shading} denotes attenuation ($|J|$ decreases). The same warming can either strengthen or weaken the bias depending solely on the curvature and temperature dependence of the underlying trait--performance relationship.}
  \label{fig:conceptual}
\end{figure}

\begin{figure}[p]
  \centering
  \includegraphics[width=\textwidth]{/home/urtzi/Documents/papers/individuals_DEB/figures/2_final/fig2-jensen-female.pdf}
  \caption{Jensen bias ($J$) between population-average and mean individual-level predictions for females under the current climate. Negative values indicate underprediction by the population-average parameterisation, whereas positive values indicate overprediction; values close to zero indicate close agreement. The colour scale is symmetric and centred at zero. The black outline delimits the current distribution of \emph{I.~alpestris} \parencite{iucn2021}. $T_b$, body temperature.}
  \label{fig:map-jensen}
\end{figure}

\begin{figure}[p]
  \centering
  \includegraphics[width=\textwidth]{/home/urtzi/Documents/papers/individuals_DEB/figures/2_final/fig3-warming-female.pdf}
   \caption{Change in Jensen bias under \qty{+4}{\celsius} warming for females. Each map displays the change in absolute Jensen bias ($\Delta|J|$, left, with colour bar) for each performance metric. The curve (descriptive GAM smoother) at the right of each map shows how $\Delta|J|$ changes along the mean annual temperature ($\overline{T}_a$) gradient. Positive values (\textit{yellow}) indicate that warming increases the disagreement between population-average and mean individual-level predictions (amplification), whereas negative values (\textit{purple}) indicate reduced disagreement (attenuation). The black outline delimits the current distribution of \emph{I.~alpestris} \parencite{iucn2021}. $T_b$, body temperature.}
  \label{fig:map-jensen-warming}
\end{figure}

\begin{figure}[p]
  \centering
  \includegraphics[width=\textwidth]{/home/urtzi/Documents/papers/individuals_DEB/figures/2_final/fig4-species-vs-ind-female.pdf}
  \caption{Population-average and mean individual-level predictions and projected changes of lifespan, maximum body mass and lifetime fecundity for females under the current climate and \qty{+2}{\celsius} and \qty{+4}{\celsius} warming. For each trait, the top row shows absolute predictions from the population-average parameterisation and the mean of the individual-level parameterisations under each scenario, on a common trait-specific scale (\textit{green}). The bottom row shows the projected change under warming relative to the corresponding current-climate prediction ($P_\mathrm{future} - P_\mathrm{current}$), normalised within each trait by the maximum absolute change across all grid cells, parameterisations and scenarios (scale from $-1$ to 1). Negative values (\textit{red}) indicate decreases under warming, positive values (\textit{blue}) indicate increases. The black outline delimits the current distribution of \emph{I.~alpestris} \parencite{iucn2021}.}
  \label{fig:map-suitability}
\end{figure}







\end{document}
